% Part: computability
% Chapter: computability-theory
% Section: equiv-ce-defs

\documentclass[../../../include/open-logic-section]{subfiles}

\begin{document}

% SEGMENT OLP-0239-S01

\olfileid{cmp}{thy}{eqc}

\olsection[க. எ. கணங்களின் வரையறைகள்]{கணிக்கத்தக்கவாறு
  எண்ணிடத்தக்க கணங்களின் சமான வரையறைகள்}


கணிக்கத்தக்கவாறு எண்ணிடத்தக்கது என்பதன் பொருளைக் கூறும் முக்கியமான
பல சமான கூற்றுகளைப் பின்வரும் தேற்றம் வழங்குகிறது.

\begin{thm}
\ollabel{thm:ce-equiv}
$S$ இயல் எண்களின் ஒரு கணம் ஆகட்டும். அப்போது பின்வருவன சமானவை:
\begin{enumerate}
\item\ollabel{case:ce} $S$ கணிக்கத்தக்கவாறு எண்ணிடத்தக்கது.
\item\ollabel{case:ran-pc} $S$ ஒரு \emph{பகுதி} கணிக்கத்தக்க சார்பின் வீச்சு.
\item\ollabel{case:ran-prim} $S$ வெற்றுக்கணம், அல்லது ஒரு முதன்மை
  மீள்வரையறைச் சார்பின் வீச்சு.
\item\ollabel{case:ce-domain} $S$ ஒரு பகுதி கணிக்கத்தக்க சார்பின்
  \emph{வரையறைப்புலம்}.
\end{enumerate}
\end{thm}

\begin{explain}
முதல் மூன்று கூறுகளும், வெற்றல்லாத எந்தக் கணிக்கத்தக்கவாறு
எண்ணிடத்தக்க கணத்தையும் ஒரு கணிக்கத்தக்க சார்பாலோ, பகுதி
கணிக்கத்தக்க சார்பாலோ, முதன்மை மீள்வரையறைச் சார்பாலோ எண்ணிடலாம்
என்று சமானமாக எடுத்துக்கொள்ளலாம் என்கின்றன. நான்காவது கூறு,
$S$ கணிக்கத்தக்கவாறு எண்ணிடத்தக்கது எனில், ஏதோ ஒரு சுட்டெண்~$e$ க்கு
\[
S = \Setabs{x}{\cfind{e}(x) \fdefined}.
\]
என்று கூறுகிறது. வேறுவிதமாகக் கூறினால், $\cfind{e}$ இன் கணிப்பு
எந்த உள்ளீடுகளில் நிற்கிறதோ, அவற்றின் கணமே~$S$. இதனால்,
கணிக்கத்தக்கவாறு எண்ணிடத்தக்க கணங்கள் சில நேரங்களில்
\emph{அரைத் தீர்மானிக்கத்தக்கவை} என்று அழைக்கப்படுகின்றன:
ஒரு எண் கணத்தில் இருந்தால் இறுதியில் ``ஆம்,'' என்ற விடை கிடைக்கும்;
அது இல்லாவிட்டால் ஒருபோதும் ``இல்லை'' என்ற விடை கிடைக்காது!{}
\end{explain}

\begin{proof}
ஒவ்வொரு முதன்மை மீள்வரையறைச் சார்பும் கணிக்கத்தக்கது; ஒவ்வொரு
கணிக்கத்தக்க சார்பும் பகுதி கணிக்கத்தக்கது. எனவே
\olref{case:ran-prim} ஆனது \olref{case:ce} ஐ உட்கிடையாக்குகிறது;
\olref{case:ce} ஆனது~\olref{case:ran-pc} ஐ உட்கிடையாக்குகிறது.
($S$ வெற்றுக்கணம் எனில், எங்குமே வரையறுக்கப்படாத பகுதி கணிக்கத்தக்க
சார்பின் வீச்சாக~$S$ உள்ளது என்பதைக் கவனிக்க.) \olref{case:ran-pc}
ஆனது \olref{case:ran-prim} ஐ உட்கிடையாக்குகிறது என்று காட்டினால்,
முதல் மூன்று கூறுகளும் சமானவை என்று காட்டியிருப்போம்.

ஆகவே, பகுதி கணிக்கத்தக்க சார்பு $\cfind{e}$ இன் வீச்சாக~$S$
இருக்கிறது எனக் கொள்க. $S$ வெற்றுக்கணம் எனில் நிரூபணம் முடிந்தது.
இல்லையெனில், $a$ ஆனது~$S$ இன் ஏதாவது ஓர் உறுப்பு ஆகட்டும்.
க்ளீனியின் இயல்பு வடிவத் தேற்றத்தால்,
\[
\cfind{e}(x) = U(\umin{s}{T(e, x, s)}).
\]
என்று எழுதலாம். குறிப்பாக, ஏதாவது~$s$ க்கு $T(e, x, s)$ உம்
$U(s) = y$ உம் இருந்தால், எனில் மட்டுமே,
$\cfind{e}(x) \fdefined$ மற்றும் $= y$. $f(z)$ ஐப் பின்வருமாறு வரையறுக்க:
\[
f(z) = \begin{cases}
  U((z)_1) & \text{$T(e, (z)_0, (z)_1)$ எனில்} \\
  a        & \text{இல்லையெனில்.}
\end{cases}
\]
$T$ உம் $U$ உம் முதன்மை மீள்வரையறைச் சார்புகளாக இருப்பதால்,
$f$ உம் அவ்வாறே உள்ளது. \iftag{TMs}{டியூரிங் பொறிகளின் மொழியில்:
$z$ ஆனது $\tuple{(z)_0, (z)_1}$ என்ற சோடியை குறிமுறையாக்குகிறது என்றும்,
$(z)_1$ ஆனது உள்ளீடு $(z)_0$ இல் பொறி~$M_e$ இன் நிற்கும் கணிப்பு
என்றும் கொள்க. அப்போது $f$ அந்தக் கணிப்பின் வெளியீட்டைத் தருகிறது;
இல்லையெனில் அது~$a$ ஐத் தருகிறது.}

$S$ ஆனது~$f$ இன் வீச்சு என்று, அதாவது ஒவ்வோர் இயல் எண்~$y$ க்கும்
$y \in S$ எனில், எனில் மட்டுமே, அது~$f$ இன் வீச்சில் இருக்கும்
என்று நாம் காட்ட வேண்டும். முன்னோக்குத் திசையில் $y \in S$ எனக்
கொள்க. அப்போது $y$ ஆனது $\cfind{e}$ இன் வீச்சில் இருப்பதால்,
ஏதோ $x$ மற்றும்~$s$ க்கு $T(e,x,s)$ மெய்யாகவும் $U(s) = y$
ஆகவும் இருக்கும். ஆனால் அப்போது $y = f(\tuple{x,s})$. மறுதிசையில்,
$y$ ஆனது~$f$ இன் வீச்சில் உள்ளது எனக் கொள்க. அப்போது $y = a$,
அல்லது ஏதோ~$z$ க்கு $T(e,(z)_0,(z)_1)$ உம்
$U((z)_1) = y$ உம் இருக்கும். பிந்தைய நேர்வில்
$\cfind{e}(x) \fdefined = y$ என்பதால், இரு நேர்வுகளிலும் $y$ ஆனது~$S$ இல் உள்ளது.

($\cfind{e}(x) \fdefined = y$ என்ற குறியீடு
``$\cfind{e}(x)$ வரையறுக்கப்பட்டு $y$ க்குச் சமம்'' என்று
பொருள்படும். $\cfind{e}(x) = y$ என்றே பயன்படுத்தியிருக்கலாம்;
ஆனால் நாம் ஒரு பகுதிச் சார்பைக் கையாள்கிறோம் என்பதை நினைவூட்டக்
கூடுதல் அம்புக்குறி சில நேரங்களில் உதவுகிறது.)

\olref{thm:ce-equiv} இன் நிரூபணத்தை முடிக்க, \olref{case:ce} உம்
~\olref{case:ce-domain} உம் சமானவை என்று காட்டினால் போதும். முதலில்
\olref{case:ce} ஆனது~\olref{case:ce-domain} ஐ உட்கிடையாக்குகிறது
என்று காட்டுவோம். கணிக்கத்தக்க சார்பு~$f$ ஆல் $S$ எண்ணிடப்படுகிறது,
அதாவது
\[
S = \Setabs{y}{\text{ஏதோ $x$ க்கு, } f(x) = y}.
\]
எனக் கொள்க. இப்போது
\[
g(y) = \umin{x}{(f(x) = y)}.
\]
என்க. அப்போது $g$ ஒரு பகுதி கணிக்கத்தக்க சார்பு; மேலும்
ஏதோ~$x$ க்கு $f(x) = y$ எனில், எனில் மட்டுமே, $g(y)$
வரையறுக்கப்பட்டுள்ளது. வேறுவிதமாகக் கூறினால், $g$ இன் வரையறைப்புலம்
~$f$ இன் வீச்சாகும். \iftag{TMs}{டியூரிங் பொறிகளின் மொழியில்:
~$S$ இன் உறுப்புகளை எண்ணிடும் டியூரிங் பொறி~$F$ கொடுக்கப்பட்டால்,
ஒரு கொடுக்கப்பட்ட உறுப்பு கணத்தில் உள்ளதா என்று அறிய~$F$ இன்
வெளியீடுகளில் தேடி, இருந்தால் நின்று, இல்லாவிட்டால் என்றென்றும்
தேடிக்கொண்டிருப்பதன் மூலம்~$S$ ஐ அரைத் தீர்மானிக்கும் டியூரிங்
பொறியாக~$G$ ஐ எடுத்துக்கொள்க.}{}

இறுதியாக, \olref{case:ce-domain} ஆனது~\olref{case:ce} ஐ
உட்கிடையாக்குகிறது என்று காட்ட, பகுதி கணிக்கத்தக்க சார்பு
~$\cfind{e}$ இன் வரையறைப்புலமாக~$S$ உள்ளது, அதாவது
\[
S = \Setabs{x}{\cfind{e}(x) \fdefined}.
\]
எனக் கொள்க. $S$ வெற்றுக்கணம் எனில் நிரூபணம் முடிந்தது;
இல்லையெனில் $a$ ஆனது~$S$ இன் ஏதாவது ஓர் உறுப்பு ஆகட்டும்.
$f$ ஐப் பின்வருமாறு வரையறுக்க:
\[
f(z) = \begin{cases}
(z)_0 & \text{$T(e,(z)_0,(z)_1)$ எனில்} \\
a & \text{இல்லையெனில்.}
\end{cases}
\]
அப்போது மேலே கண்டபடி, $\cfind{e}(x) \fdefined$ எனில், எனில்
மட்டுமே, அதாவது $x \in S$ எனில், எனில் மட்டுமே, எண்~$x$ ஆனது
~$f$ இன் வீச்சில் உள்ளது. \iftag{TMs}{டியூரிங் பொறிகளின் மொழியில்:
~$S$ ஐ அரைத் தீர்மானிக்கும் பொறி $M_e$ கொடுக்கப்பட்டால், சாத்தியமான
எல்லா டியூரிங் பொறிக் கணிப்புகளையும் ஒன்றன்பின் ஒன்றாக ஆராய்ந்து,
நிற்கும் கணிப்புகளுக்குரிய உள்ளீடுகளைத் தருவதன் மூலம்~$S$ இன்
உறுப்புகளை எண்ணிடுக.}{}
\end{proof}

கணிக்கத்தக்கவாறு எண்ணிடத்தக்க கணங்களை எண்ணிடுவதற்கான வசதியான
முறையை, \olref{case:ce-domain} என்ற கூறு~\olref{thm:ce-equiv} இல்
வழங்குகிறது:
ஒவ்வொரு~$e$ க்கும், $\cfind{e}$ இன் வரையறைப்புலத்தை $W_e$ குறிக்கட்டும்,
அதாவது
\[
W_e = \Setabs{x}{\cfind{e}(x) \fdefined}.
\]
அப்போது $A$ எந்தக் கணிக்கத்தக்கவாறு எண்ணிடத்தக்க கணமாக இருந்தாலும்,
ஏதோ~$e$ க்கு $A = W_e$.

கணிக்கத்தக்கவாறு எண்ணிடத்தக்க கணங்களுக்கான மற்றொரு இயல்பாக்கத்தைப்
பின்வரும் தேற்றம் வழங்குகிறது.

\begin{thm}
\ollabel{thm:exists-char}
கணிக்கத்தக்க உறவு $R(x,y)$ ஒன்று இருந்து
\[
S = \Setabs{ x }{ \lexists[y][R(x,y)] }.
\]
எனில், எனில் மட்டுமே, கணம் $S$ கணிக்கத்தக்கவாறு எண்ணிடத்தக்கது.
\end{thm}

\begin{proof}
முன்னோக்குத் திசையில், $S$ கணிக்கத்தக்கவாறு எண்ணிடத்தக்கது எனக்
கொள்க. அப்போது ஏதோ $e$ க்கு $S = W_e$. இந்த~$e$ மதிப்புக்கு,
$S$ ஐ
\[
S = \Setabs{ x }{ \lexists[y][T(e, x, y)] }.
\]
என்று எழுதலாம். மறுதிசையில்,
$S = \Setabs{ x }{ \lexists[y][R(x, y)] }$ எனக் கொள்க.
$f$ ஐப் பின்வருமாறு வரையறுக்க:
\[
f(x) \simeq \umin{y}{R(x, y)}.
\]
அப்போது $f$ ஒரு பகுதி கணிக்கத்தக்க சார்பு; $S$ ஆனது~$f$ இன்
வரையறைப்புலம்.
\end{proof}


\end{document}
